\documentclass[12pt]{article}
\usepackage[utf8]{inputenc}
\usepackage[T1]{fontenc}
\usepackage[margin=1in]{geometry}
\usepackage{setspace}
\usepackage{xcolor}
\usepackage[colorlinks=true, urlcolor=blue!60!black, linkcolor=blue!60!black]{hyperref}
\setlength{\parindent}{0em}
\setlength{\parskip}{0.6em}
\renewcommand{\baselinestretch}{1.0}
\begin{document}
\begin{center}
{\LARGE \textit{Zelman v. Simmons-Harris}, 536 U.S. 639 (2002)\par}
\vspace{0.5em}
{\large Case Brief --- Two-Page Edition\par}
\end{center}

\section{Memory Jogger}
Cleveland's voucher program survived 5--4: neutral aid that reaches religious schools only through families' genuine and independent private choice does not violate the Establishment Clause (\href{https://supreme.justia.com/cases/federal/us/536/639/}{Justia}).

\section{Detailed Case Facts}
After Cleveland's failing schools were placed under state control, Ohio offered low-income families tuition scholarships usable at private schools, religious or secular, or adjacent public schools, plus tutoring grants for public-school students. Of 56 participating private schools, 46 were religious, and 96\% of scholarship students attended religious schools. Cleveland also offered community and magnet schools.

\section{Procedural History}
Taxpayers won summary judgment, and a divided Sixth Circuit affirmed, finding that the program advanced religion.

\section{Judicial Votes}
Reversed, 5--4. Rehnquist wrote for O'Connor, Scalia, Kennedy, and Thomas; O'Connor and Thomas concurred. Stevens, Souter (with Stevens, Ginsburg, and Breyer), and Breyer (with Stevens and Souter) dissented.

\section{Holding}
A religion-neutral program aiding a broad class of citizens, who direct the aid to religious schools wholly by their own genuine and independent choice, does not violate the Establishment Clause.

\section{Analysis of Opinions}
\textbf{Majority:} Under \textit{Mueller}, \textit{Witters}, and \textit{Zobrest}, true private choice breaks the circuit between government and religion. The program was neutral and even disfavored religious schools financially, and genuine choice must be measured across all of Cleveland's options, including community and magnet schools. \textbf{O'Connor:} indirect aid like Pell Grants is longstanding. \textbf{Thomas:} school choice emancipates urban minority children. \textbf{Souter:} the Court traded \textit{Everson}'s no-aid rule for verbal formalism, leaving parents a Hobson's choice. \textbf{Stevens and Breyer:} public funding of religious teaching invites strife. \emph{Later:} \textit{Espinoza} and \textit{Carson v. Makin} made inclusion of religious schools mandatory once a State funds private schooling.

\section{Examples: Future Applications}
\textbf{Same-Side:} savings accounts that flow entirely to religious academies in a county lacking secular private schools are valid. \textbf{Opposite-Side:} vouchers redeemable only at one denomination's schools are not neutral. \textbf{Fence-Sitter:} offenders choosing among mostly faith-based rehabilitation programs while the secular one is waitlisted.

\section{Critique}
The result turned on the denominator: counting magnet and community schools cut religious enrollment from 96\% to under 20\%. With \textit{Carson}, States must fund only public schools or fund religious and secular private schools alike.

\section{Key Quotations}
\begin{quote}
``[W]here a government aid program is neutral with respect to religion, and provides assistance directly to a broad class of citizens who, in turn, direct government aid to religious schools wholly as a result of their own genuine and independent private choice, the program is not readily subject to challenge under the Establishment Clause.''
\end{quote}
\noindent --- Chief Justice Rehnquist, for the Court, 536 U.S. at 652.

\end{document}
